\begin{minipage}{\linewidth}
\vspace{1em}\noindent\textbf{1. Регулировка зеркал при перевозке пассажира. Пассажир должен:}

\begin{itemize}[nosep,leftmargin=2em]
\item [A.] Стоять позади мотоцикла.
\item [B.] Сесть на сиденье рядом с вами.
\item [C.] Регулировать зеркало, пока вы сидите на мотоцикле.
\item [D.] Стоять рядом с мотоциклом в мёртвой зоне.

\end{itemize}
\end{minipage}


\begin{minipage}{\linewidth}
\vspace{1em}\noindent\textbf{2. Пассажир должен:}

\begin{itemize}[nosep,leftmargin=2em]
\item [A.] Держаться за ремень сиденья.
\item [B.] Держаться за боковую часть сиденья.
\item [C.] Держаться за талию водителя.
\item [D.] Не ездить без спинки.

\end{itemize}
\end{minipage}


\begin{minipage}{\linewidth}
\vspace{1em}\noindent\textbf{3. Когда мотоцикл останавливается на перекрёстке, пассажир должен:}

\begin{itemize}[nosep,leftmargin=2em]
\item [A.] Держать ноги на подножках.
\item [B.] Поставить одну ногу на землю.
\item [C.] Поставить обе ноги на землю.
\item [D.] Отодвинуться назад на сиденье.

\end{itemize}
\end{minipage}


\begin{minipage}{\linewidth}
\vspace{1em}\noindent\textbf{4. При езде с пассажиром:}

\begin{itemize}[nosep,leftmargin=2em]
\item [A.] Вы будете лучше удерживать равновесие.
\item [B.] Вы сможете быстрее замедлиться.
\item [C.] Разгон и торможение займут больше времени.
\item [D.] У вас будет меньше шансов попасть в занос.

\end{itemize}
\end{minipage}


\begin{minipage}{\linewidth}
\vspace{1em}\noindent\textbf{5. При езде с пассажирами вам следует:}

\begin{itemize}[nosep,leftmargin=2em]
\item [A.] Начинать тормозить раньше чем обычно.
\item [B.] Использовать для торможения только передний тормоз.
\item [C.] Никогда не разговаривать с ним.
\item [D.] Начинать движение в центре полосы движения.

\end{itemize}
\end{minipage}


\begin{minipage}{\linewidth}
\vspace{1em}\noindent\textbf{6. При перевозке грузов:}

\begin{itemize}[nosep,leftmargin=2em]
\item [A.] Располагайте груз позади задней оси.
\item [B.] Закрепите груз с помощью веревки.
\item [C.] Закрепите груз к раме сзади сиденья.
\item [D.] Располагайте груз низко и как можно ближе к передней части.

\end{itemize}
\end{minipage}


\begin{minipage}{\linewidth}
\vspace{1em}\noindent\textbf{7. Группа байкеров переходит из шахматного порядка в колонну по одному, в случае:}

\begin{itemize}[nosep,leftmargin=2em]
\item [A.] Езды ночью.
\item [B.] При интенсивном дорожном движении.
\item [C.] При остановке на перекрестках.
\item [D.] При прохождении поворотов.

\end{itemize}
\end{minipage}


\begin{minipage}{\linewidth}
\vspace{1em}\noindent\textbf{8. На этом изображении мотоциклисты едут в шахматном порядке. Дистанция между A и C должна составлять не менее:}

\vspace{1ex}
\centering
\includegraphics[width=0.8\linewidth]{ q_8.jpeg }
\captionof*{figure}{ Четыре мотоцикла (обозначенные буквами A, B, C, D) движутся в шахматном порядке по многополосной дороге: A — лидер в левой части, B — следует за лидером в правой части, C — третий, в левой части и позади лидера. }

\begin{itemize}[nosep,leftmargin=2em]
\item [A.] Четыре секунды.
\item [B.] Длина одного автомобиля.
\item [C.] Одна секунда.
\item [D.] Две секунды.

\end{itemize}
\end{minipage}


\begin{minipage}{\linewidth}
\vspace{1em}\noindent\textbf{9. Мотоциклисты, едущие в шахматном порядке, будут обгонять автомобиль. После того как лидер закончит обгон, он должен:}

\begin{itemize}[nosep,leftmargin=2em]
\item [A.] Занять правую часть полосы движения.
\item [B.] Занять центральную часть полосы движения.
\item [C.] Занять левую часть полосы движения.
\item [D.] Подать сигнал остальным, чтобы они следовали за ним.

\end{itemize}
\end{minipage}


\begin{minipage}{\linewidth}
\vspace{1em}\noindent\textbf{10. Алкоголь попадает в кровоток и быстро достигает мозга. На что это влияет в первую очередь?}

\begin{itemize}[nosep,leftmargin=2em]
\item [A.] Способность принимать адекватные решения.
\item [B.] Речь.
\item [C.] Чувство равновесия.
\item [D.] Зрение.

\end{itemize}
\end{minipage}


\begin{minipage}{\linewidth}
\vspace{1em}\noindent\textbf{11. Алкоголь в первую очередь влияет на:}

\begin{itemize}[nosep,leftmargin=2em]
\item [A.] Зрение.
\item [B.] Речь.
\item [C.] Способность принимать адекватные решения.
\item [D.] Чувство равновесия.

\end{itemize}
\end{minipage}


\begin{minipage}{\linewidth}
\vspace{1em}\noindent\textbf{12. Примерно половина всех смертельных случаев в результате ДТП с участием мотоциклов происходит по причине:}

\begin{itemize}[nosep,leftmargin=2em]
\item [A.] Состоянием дорог.
\item [B.] Шлемом мотоциклиста.
\item [C.] Из-за воздействия алкоголя.
\item [D.] Механическая неисправность.

\end{itemize}
\end{minipage}


\begin{minipage}{\linewidth}
\vspace{1em}\noindent\textbf{13. Если вы не контролируете количество выпитого:}

\begin{itemize}[nosep,leftmargin=2em]
\item [A.] Езжайте медленнее.
\item [B.] Перед поездкой выдержите как минимум час после каждой порции алкоголя.
\item [C.] Занимайтесь спортом или танцуйте, чтобы вывести алкоголь из организма.
\item [D.] Пейте горячий кофе, чтобы оставаться бодрым

\end{itemize}
\end{minipage}


\begin{minipage}{\linewidth}
\vspace{1em}\noindent\textbf{14. Исследования показывают, что почти половина всех мотоциклистов, погибших в ДТП:}

\begin{itemize}[nosep,leftmargin=2em]
\item [A.] Имели плохое зрение.
\item [B.] Употребляли алкоголь.
\item [C.] Превышали скорость.
\item [D.] Не обладали достаточным опытом

\end{itemize}
\end{minipage}


\begin{minipage}{\linewidth}
\vspace{1em}\noindent\textbf{15. Какое количество алкоголя начинает влять на способность ясно мыслить и безопасно управлять мотоциклом?}

\begin{itemize}[nosep,leftmargin=2em]
\item [A.] 1 порция алкоголя.
\item [B.] 6 порций алкоголя.
\item [C.] 3 порции алкоголя.
\item [D.] 2 порции алкоголя.

\end{itemize}
\end{minipage}


\begin{minipage}{\linewidth}
\vspace{1em}\noindent\textbf{16. Если вы выпьете 4 стандартных порции алкоголя, какое минимальное количество алкоголя останется в вашем организме через 2 часа?}

\begin{itemize}[nosep,leftmargin=2em]
\item [A.] 2 порции.
\item [B.] 3 порции.
\item [C.] 1 порция.
\item [D.] 4 порции.

\end{itemize}
\end{minipage}


\begin{minipage}{\linewidth}
\vspace{1em}\noindent\textbf{17. Для максимальной безопасности ваш шлем должен:}

\begin{itemize}[nosep,leftmargin=2em]
\item [A.] Плотно прилегать у основания; свободно — в верхней части.
\item [B.] Сидеть свободно по всему периметру.
\item [C.] Плотно прилегать по всему периметру.
\item [D.] Носить только во время длительных поездок.

\end{itemize}
\end{minipage}


\begin{minipage}{\linewidth}
\vspace{1em}\noindent\textbf{18. Самое важное, что вы можете сделать, чтобы повысить свои шансы на выживание при аварии это:}

\begin{itemize}[nosep,leftmargin=2em]
\item [A.] Опрокинуть мотоцикл до столкновения с препятствием.
\item [B.] Носить надежно закрепленный шлем.
\item [C.] Хорошо знать свой мотоцикл.
\item [D.] Ездить по центральной части полосы движения.

\end{itemize}
\end{minipage}


\begin{minipage}{\linewidth}
\vspace{1em}\noindent\textbf{19. Ношение шлема может снизить вероятность получения смертельной травмы головы при:}

\begin{itemize}[nosep,leftmargin=2em]
\item [A.] Только при столкновении на низкой скорости.
\item [B.] Только при ДТП с участием одного транспортного средства.
\item [C.] Любом ДТП, независимо от скорости.
\item [D.] Только при столкновениях на высокой скорости.

\end{itemize}
\end{minipage}


\begin{minipage}{\linewidth}
\vspace{1em}\noindent\textbf{20. Ношение шлема:}

\begin{itemize}[nosep,leftmargin=2em]
\item [A.] Снизит вероятность травм шеи и головы.
\item [B.] Увеличит вероятность травмы шеи.
\item [C.] Ухудшит обзор на две трети.
\item [D.] Ни уменьшит, ни увеличит вероятность травмы головы.

\end{itemize}
\end{minipage}


\begin{minipage}{\linewidth}
\vspace{1em}\noindent\textbf{21. Что из перечисленного ниже обеспечивает наилучшую защиту глаз и лица во время езды?}

\begin{itemize}[nosep,leftmargin=2em]
\item [A.] Защитные очки.
\item [B.] Ветровое стекло (щиток на мотоцикле).
\item [C.] Защитный щиток шлема.
\item [D.] Солнцезащитные очки.

\end{itemize}
\end{minipage}


\begin{minipage}{\linewidth}
\vspace{1em}\noindent\textbf{22. При езде на незнакомом мотоцикле:}

\begin{itemize}[nosep,leftmargin=2em]
\item [A.] Не переключайтесь выше третьей передачи.
\item [B.] Держитесь правой части полосы движения.
\item [C.] Попросите владельца поехать с вами.
\item [D.] Оставляйте больше дистанции для торможения.

\end{itemize}
\end{minipage}


\begin{minipage}{\linewidth}
\vspace{1em}\noindent\textbf{23. Аварийный выключатель двигателя обычно расположен:}

\begin{itemize}[nosep,leftmargin=2em]
\item [A.] Рядом с ручагом сцепления.
\item [B.] На левой ручке руля.
\item [C.] Рядом с правой ручкой
\item [D.] На верхней части топливного бака

\end{itemize}
\end{minipage}


\begin{minipage}{\linewidth}
\vspace{1em}\noindent\textbf{24. При правильной регулировке зеркала вы можете видеть:}

\begin{itemize}[nosep,leftmargin=2em]
\item [A.] Полосу движения рядом с вами.
\item [B.] Полосу движения позади вас.
\item [C.] Часть полосы движения рядом с вами и позади вас.
\item [D.] Всё, что находится сзади и сбоку.

\end{itemize}
\end{minipage}


\begin{minipage}{\linewidth}
\vspace{1em}\noindent\textbf{25. При замедлении перед обычной остановкой:}

\begin{itemize}[nosep,leftmargin=2em]
\item [A.] Используйте только задний тормоз.
\item [B.] Уменьшите подачу газа и используйте только передний тормоз.
\item [C.] При замедлении переключайтесь на более низкую передачу.
\item [D.] Выжмите сцепление и двигайтесь накатом.

\end{itemize}
\end{minipage}


\begin{minipage}{\linewidth}
\vspace{1em}\noindent\textbf{26. При повороте следует:}

\begin{itemize}[nosep,leftmargin=2em]
\item [A.] Следить за центральной линией.
\item [B.] Смотреть прямо вперед.
\item [C.] Смотреть на ограждение.
\item [D.] Смотреть вперед, в сторону поворота.

\end{itemize}
\end{minipage}


\begin{minipage}{\linewidth}
\vspace{1em}\noindent\textbf{27. Какое из следующих утверждений верно:}

\begin{itemize}[nosep,leftmargin=2em]
\item [A.] Нажмите на рукоять руля слева, наклонитесь влево, поедитете вправо.
\item [B.] Нажми с левой стороны, наклонись влево, поедешь влево.
\item [C.] Нажми с левой стороны, наклонись вправо, поедешь вправо.
\item [D.] Нажми с правой стороны, наклонись вправо, поедешь влево.

\end{itemize}
\end{minipage}


\begin{minipage}{\linewidth}
\vspace{1em}\noindent\textbf{28. Во время поворота лучше всего:}

\begin{itemize}[nosep,leftmargin=2em]
\item [A.] Сохранять постоянную скорость или плавно ускоряться.
\item [B.] Снизить скорость.
\item [C.] Использовать тормоза.
\item [D.] Ускоряться и смотреть прямо вперед.

\end{itemize}
\end{minipage}


\begin{minipage}{\linewidth}
\vspace{1em}\noindent\textbf{29. Если вам нужно замедлиться во время поворота:}

\begin{itemize}[nosep,leftmargin=2em]
\item [A.] Задействуйте задний тормоз и наклонитесь сильнее в сторону поворота.
\item [B.] Ни в коем случае не тормозите.
\item [C.] Аккуратно задействуйте оба тормоза.
\item [D.] Задействуйте только задний тормоз.

\end{itemize}
\end{minipage}


\begin{minipage}{\linewidth}
\vspace{1em}\noindent\textbf{30. Для обычной остановки используйте:}

\begin{itemize}[nosep,leftmargin=2em]
\item [A.] Сначала задний тормоз.
\item [B.] Оба тормоза и переключение на пониженную передачу.
\item [C.] Передний тормоз и переключение на пониженную передачу.
\item [D.] Только задний тормоз.

\end{itemize}
\end{minipage}


\begin{minipage}{\linewidth}
\vspace{1em}\noindent\textbf{31. У вашего мотоцикла два тормоза. Используйте оба тормоза:}

\begin{itemize}[nosep,leftmargin=2em]
\item [A.] Только при экстренной остановке.
\item [B.] При резком маневре.
\item [C.] Каждый раз, когда вы замедляетесь или останавливаетесь.
\item [D.] Только при обычной остановке.

\end{itemize}
\end{minipage}


\begin{minipage}{\linewidth}
\vspace{1em}\noindent\textbf{32. При обычном повороте:}

\begin{itemize}[nosep,leftmargin=2em]
\item [A.] Наклоняйте мотоцикл сильнее, чем тело.
\item [B.] Наклоняйте мотоцикл и тело под одинаковым углом.
\item [C.] Не наклоняйтесь.
\item [D.] Наклоняйте тело сильнее, чем мотоцикл.

\end{itemize}
\end{minipage}


\begin{minipage}{\linewidth}
\vspace{1em}\noindent\textbf{33. При медленных крутых поворотах:}

\begin{itemize}[nosep,leftmargin=2em]
\item [A.] Наклоняйте только мотоцикл, а тело держите прямо.
\item [B.] Наклоняйте тело сильнее, чем мотоцикл.
\item [C.] Наклоняйтесь вместе с мотоциклом.
\item [D.] Не наклоняйтесь.

\end{itemize}
\end{minipage}


\begin{minipage}{\linewidth}
\vspace{1em}\noindent\textbf{34. Чтобы снизить скорость перед поворотом:}

\begin{itemize}[nosep,leftmargin=2em]
\item [A.] Переключитесь на нейтраль.
\item [B.] Зажмите сцепление и задействуйте оба тормоза.
\item [C.] Используйте только задний тормоз.
\item [D.] Закройте дроссель и задействуйте оба тормоза.

\end{itemize}
\end{minipage}


\begin{minipage}{\linewidth}
\vspace{1em}\noindent\textbf{35. При повороте важно:}

\begin{itemize}[nosep,leftmargin=2em]
\item [A.] Смотрите прямо вперед.
\item [B.] Снизить скорость после входа в поворот.
\item [C.] Снизить скорость при выходе из поворота.
\item [D.] Смотреть сквозь поворот туда, куда вы хотите ехать.

\end{itemize}
\end{minipage}


\begin{minipage}{\linewidth}
\vspace{1em}\noindent\textbf{36. Важно искать на дороге пути отхода:}

\begin{itemize}[nosep,leftmargin=2em]
\item [A.] Через переулки.
\item [B.] На перекрестках.
\item [C.] При езде на одолженном мотоцикле.
\item [D.] На поворотах.

\end{itemize}
\end{minipage}


\begin{minipage}{\linewidth}
\vspace{1em}\noindent\textbf{37. Чтобы правильно принимать решения в дорожном движении, сначала нужно:}

\begin{itemize}[nosep,leftmargin=2em]
\item [A.] Уметь резко сворачивать.
\item [B.] Смотреть вперёд.
\item [C.] Следить за пешеходами.
\item [D.] Уметь быстро останавливаться.

\end{itemize}
\end{minipage}


\begin{minipage}{\linewidth}
\vspace{1em}\noindent\textbf{38. Водитель автомобиля подаёт сигнал о повороте налево. Мотоциклист должен:}

\vspace{1ex}
\centering
\includegraphics[width=0.8\linewidth]{ q_38.jpeg }
\captionof*{figure}{ Сцена на перекрёстке, где мотоцикл и автомобиль встречаются на перекрёстке. Мотоцикл движется прямо, а автомобиль поворачивает перед мотоциклом. Мотоцикл имеет преимущество в движении. }

\begin{itemize}[nosep,leftmargin=2em]
\item [A.] Сохранять скорость и держаться в центре полосы движения.
\item [B.] Снизить скорость и держаться подальше от приближающихся транспортных средств.
\item [C.] Увеличить скорость и двигаться влево.
\item [D.] Увеличить скорость и оставаться в центре полосы движения.

\end{itemize}
\end{minipage}


\begin{minipage}{\linewidth}
\vspace{1em}\noindent\textbf{39. Вам следует уделять особое внимание проверке обстановки сзади:}

\begin{itemize}[nosep,leftmargin=2em]
\item [A.] Перед входом в поворот.
\item [B.] При приближении к перекрёстку.
\item [C.] Когда везёте пассажира.
\item [D.] Перед тем, как снизить скорость.

\end{itemize}
\end{minipage}


\begin{minipage}{\linewidth}
\vspace{1em}\noindent\textbf{40. При подготовке к обгону обычно лучше:}

\begin{itemize}[nosep,leftmargin=2em]
\item [A.] Подать сигнал, посмотреть в левое зеркало,  посмотреть налево повернув голову.
\item [B.] Подать сигнал, перестроиться в левый ряд, ускориться.
\item [C.] Подать сигнал, посмотреть в левое зеркало.
\item [D.] Включить поворотник, посмотреть в левое зеркало, ускориться.

\end{itemize}
\end{minipage}


\begin{minipage}{\linewidth}
\vspace{1em}\noindent\textbf{41. Движение по центру полосы:}

\begin{itemize}[nosep,leftmargin=2em]
\item [A.] Мешает другим участникам движения увидеть вас.
\item [B.] Является самым безопасным местом при приближении к перекрёсткам.
\item [C.] Не позволяет другим транспортным средствам занимать вашу полосу.
\item [D.] По возможности следует избегать.

\end{itemize}
\end{minipage}


\begin{minipage}{\linewidth}
\vspace{1em}\noindent\textbf{42. Чтобы предотвратить попытки других водителей занять вашу полосу вместе с вами, обычно лучше всего:}

\begin{itemize}[nosep,leftmargin=2em]
\item [A.] Двигаться по левой части полосы.
\item [B.] Двигаться по центральной части полосы.
\item [C.] Ехать рядом с другим транспортным средством.
\item [D.] Двигаться по правой части полосы.

\end{itemize}
\end{minipage}


\begin{minipage}{\linewidth}
\vspace{1em}\noindent\textbf{43. Поддержание зоны безопасности (пространственного буфера) вокруг мотоцикла:}

\begin{itemize}[nosep,leftmargin=2em]
\item [A.] Необходимо только начинающим мотоциклистам.
\item [B.] Даёт вам время и пространство для реакции.
\item [C.] Позволяет лавировать в потоке при необходимости.
\item [D.] Делает вас более уязвимым.

\end{itemize}
\end{minipage}


\begin{minipage}{\linewidth}
\vspace{1em}\noindent\textbf{44. Чтобы вовремя заметить встречный транспорт и сохранить зону безопасности, как показано на иллюстрации выше:}

\vspace{1ex}
\centering
\includegraphics[width=0.8\linewidth]{ q_44.jpeg }
\captionof*{figure}{ Вид вперёд на извилистую дорогу }

\begin{itemize}[nosep,leftmargin=2em]
\item [A.] Двигайтесь по центру полосы.
\item [B.] Сместитесь вправо и отклонитесь в сторону, противоположную повороту.
\item [C.] Не наклоняйте мотоцикл.
\item [D.] Двигайтесь по левой части полосы.

\end{itemize}
\end{minipage}


\begin{minipage}{\linewidth}
\vspace{1em}\noindent\textbf{45. Вы планируете повернуть направо после остановки на перекрёстке. Обычно лучше всего:}

\begin{itemize}[nosep,leftmargin=2em]
\item [A.] Позволить авто сзади, объехать вас.
\item [B.] Выполнять поворот из центральной части полосы.
\item [C.] Выполнять поворот из правой части полосы.
\item [D.] Выполнять поворот из левой части полосы.

\end{itemize}
\end{minipage}


\begin{minipage}{\linewidth}
\vspace{1em}\noindent\textbf{46. Движение на близкой дистанции за впереди идущим транспортом:}

\begin{itemize}[nosep,leftmargin=2em]
\item [A.] Помогает видеть опасности на дороге.
\item [B.] Безопаснее в ночное время.
\item [C.] Защищает вас при плотном движении.
\item [D.] Является одной из причин ДТП по вине мотоциклистов.

\end{itemize}
\end{minipage}


\begin{minipage}{\linewidth}
\vspace{1em}\noindent\textbf{47. Двигаясь за автомобилем, вам следует стараться:}

\begin{itemize}[nosep,leftmargin=2em]
\item [A.] Ехать так, чтобы водитель видел вас в зеркало заднего вида.
\item [B.] Ехать рядом с другим автомобилем, чтобы быть заметнее.
\item [C.] Двигаться по левой стороне полосы.
\item [D.] Двигаться по правой стороне полосы.

\end{itemize}
\end{minipage}


\begin{minipage}{\linewidth}
\vspace{1em}\noindent\textbf{48. Самой частой причиной ДТП с участием мотоциклов и автомобилей на перекрёстках является:}

\begin{itemize}[nosep,leftmargin=2em]
\item [A.] Нарушение водителями преимущественного права проезда мотоциклиста.
\item [B.] То, что мотоциклистов не замечают выезжающие со второстепенных улиц автомобили.
\item [C.] Несоблюдение водителями дистанции позади мотоциклиста.
\item [D.] Непредоставление мотоциклистами преимущества встречному транспорту.

\end{itemize}
\end{minipage}


\begin{minipage}{\linewidth}
\vspace{1em}\noindent\textbf{49. ДТП с участием автомобилей и мотоциклов чаще всего происходят:}

\begin{itemize}[nosep,leftmargin=2em]
\item [A.] На перекрёстках.
\item [B.] Ночью.
\item [C.] При превышении скорости мотоциклистом.
\item [D.] На автомагистралях.

\end{itemize}
\end{minipage}


\begin{minipage}{\linewidth}
\vspace{1em}\noindent\textbf{50. Зрительный контакт с водителем на перекрёстке:}

\begin{itemize}[nosep,leftmargin=2em]
\item [A.] Является верным признаком того, что водитель, скорее всего, уступит дорогу.
\item [B.] Важен, если водитель выезжает со второстепенной улицы.
\item [C.] Не означает, что водитель действительно вас видит.
\item [D.] Вводит водителя в заблуждение.

\end{itemize}
\end{minipage}


\begin{minipage}{\linewidth}
\vspace{1em}\noindent\textbf{51. Наибольшая опасность для мотоциклиста при проезде мимо припаркованных автомобилей — это:}

\begin{itemize}[nosep,leftmargin=2em]
\item [A.] Автомобиль, заезжающий на парковочное место задним ходом.
\item [B.] Автомобиль, выезжающий с парковки на дорогу.
\item [C.] Открывающаяся дверь автомобиля.
\item [D.] Пешеход, выходящий из-за припаркованных машин.

\end{itemize}
\end{minipage}


\begin{minipage}{\linewidth}
\vspace{1em}\noindent\textbf{52. На этом рисунке световой сигнал, который скорее всего сделает мотоцикл заметным для водителя автомобиля — это:}

\vspace{1ex}
\centering
\includegraphics[width=0.8\linewidth]{ q_52.jpeg }
\captionof*{figure}{ Мотоциклист выезжает на шоссе, когда по правой полосе едет автомобиль }

\begin{itemize}[nosep,leftmargin=2em]
\item [A.] Стоп-сигнал.
\item [B.] Задний габаритный огонь.
\item [C.] Фара.
\item [D.] Указатель поворота.

\end{itemize}
\end{minipage}


\begin{minipage}{\linewidth}
\vspace{1em}\noindent\textbf{53. При движении в потоке ночью лучший способ обнаружить неровности на дороге — это:}

\begin{itemize}[nosep,leftmargin=2em]
\item [A.] Смотреть поверх едущих впереди автомобилей.
\item [B.] Включить дальний свет фар.
\item [C.] Следить за дорожными знаками.
\item [D.] Следить за задними габаритными огнями впереди идущего автомобиля.

\end{itemize}
\end{minipage}


\begin{minipage}{\linewidth}
\vspace{1em}\noindent\textbf{54. Лучший способ сделать ваш мотоцикл заметнее для других — это:}

\begin{itemize}[nosep,leftmargin=2em]
\item [A.] Двигаться по левой части полосы.
\item [B.] Надевать белый шлем.
\item [C.] Держаться в центральной части полосы.
\item [D.] Ездить с включённой фарой.

\end{itemize}
\end{minipage}


\begin{minipage}{\linewidth}
\vspace{1em}\noindent\textbf{55. Забытый включённый указатель поворота скорее всего:}

\begin{itemize}[nosep,leftmargin=2em]
\item [A.] Спровоцирует других водителей выехать на вашу траекторию движения.
\item [B.] Помешает другим заметить ваш стоп-сигнал.
\item [C.] Снизит вероятность аварии, так как другие водители будут держаться подальше.
\item [D.] Заставит других ехать слишком близко за вами.

\end{itemize}
\end{minipage}


\begin{minipage}{\linewidth}
\vspace{1em}\noindent\textbf{56. Особенно важно проверять зеркала заднего вида:}

\begin{itemize}[nosep,leftmargin=2em]
\item [A.] В дальних поездках.
\item [B.] При проезде перекрёстка.
\item [C.] При движении за другим транспортным средством.
\item [D.] Перед замедлением или остановкой.

\end{itemize}
\end{minipage}


\begin{minipage}{\linewidth}
\vspace{1em}\noindent\textbf{57. Перед перестроением на левую полосу:}

\begin{itemize}[nosep,leftmargin=2em]
\item [A.] Проверьте левое зеркало и поверните голову влево для проверки слепой зоны.
\item [B.] Проверьте правое зеркало, затем левое.
\item [C.] Проверьте левое зеркало, затем правое.
\item [D.] Подайте звуковой сигнал и ускорьтесь.

\end{itemize}
\end{minipage}


\begin{minipage}{\linewidth}
\vspace{1em}\noindent\textbf{58. Лучшая экипировка, чтобы быть заметным ночью — это:}

\begin{itemize}[nosep,leftmargin=2em]
\item [A.] Флуоресцентная одежда.
\item [B.] Белый шлем.
\item [C.] Одежда со световозвращающими элементами.
\item [D.] Красная куртка.

\end{itemize}
\end{minipage}


\begin{minipage}{\linewidth}
\vspace{1em}\noindent\textbf{59. Чтобы лучше видеть дорогу при движении за автомобилем ночью:}

\begin{itemize}[nosep,leftmargin=2em]
\item [A.] Увеличьте дистанцию и используйте ближний свет.
\item [B.] Используйте свет фар впереди идущего автомобиля.
\item [C.] Совершите обгон при безопасности и включите дальний свет.
\item [D.] Отстаньте подальше и включите дальний свет.

\end{itemize}
\end{minipage}


\begin{minipage}{\linewidth}
\vspace{1em}\noindent\textbf{60. При езде в ночное время:}

\begin{itemize}[nosep,leftmargin=2em]
\item [A.] Всегда двигайтесь по центральной полосе.
\item [B.] Ездите с включённой аварийной сигнализацией.
\item [C.] Носите световозвращающую одежду.
\item [D.] Держитесь ближе к едущему впереди транспорту для безопасности.

\end{itemize}
\end{minipage}


\begin{minipage}{\linewidth}
\vspace{1em}\noindent\textbf{61. Если при остановке вы заблокировали переднее колесо:}

\begin{itemize}[nosep,leftmargin=2em]
\item [A.] Отпустить задний тормоз.
\item [B.] Держать передний тормоз заблокированным до полной остановки.
\item [C.] Держать заблокированными оба тормоза.
\item [D.] Отпустить и снова нажать передний тормоз.

\end{itemize}
\end{minipage}


\begin{minipage}{\linewidth}
\vspace{1em}\noindent\textbf{62. На рисунке выше при попытке экстренной остановки заблокировалось заднее колесо. Как правило, лучше всего:}

\vspace{1ex}
\centering
\includegraphics[width=0.8\linewidth]{ q_62.jpeg }
\captionof*{figure}{ Перекресток улиц в пригородном районе: слева автомобиль, справа мотоцикл, приближающиеся друг к другу. }

\begin{itemize}[nosep,leftmargin=2em]
\item [A.] Сохранять положение и отпустить задний тормоз.
\item [B.] Наклонить мотоцикл и отпустить задний тормоз.
\item [C.] Отпустить передний тормоз и держать задний тормоз заблокированным.
\item [D.] Держать заднее колесо заблокированным до полной остановки.

\end{itemize}
\end{minipage}


\begin{minipage}{\linewidth}
\vspace{1em}\noindent\textbf{63. При резком объезде препятствия важно:}

\begin{itemize}[nosep,leftmargin=2em]
\item [A.] Всегда тормозить перед объездом.
\item [B.] Разделять торможение и маневр.
\item [C.] Всегда тормозить после маневра.
\item [D.] Сворачивать в сторону движения опасности.

\end{itemize}
\end{minipage}


\begin{minipage}{\linewidth}
\vspace{1em}\noindent\textbf{64. Для проезда по металлическим решетчатым настилам моста:}

\begin{itemize}[nosep,leftmargin=2em]
\item [A.] Двигаться прямо.
\item [B.] Двигаться зигзагом.
\item [C.] Пересекать под острым углом, не меняя полосы.
\item [D.] Увеличить скорость.

\end{itemize}
\end{minipage}


\begin{minipage}{\linewidth}
\vspace{1em}\noindent\textbf{65. При переезде через предмет на дороге:}

\begin{itemize}[nosep,leftmargin=2em]
\item [A.] Оставаться в сиденье, пока не переедете его.
\item [B.] Приподняться над сиденьем.
\item [C.] Переезжать его по самому краю.
\item [D.] Сжать бензобак коленями.

\end{itemize}
\end{minipage}


\begin{minipage}{\linewidth}
\vspace{1em}\noindent\textbf{66. При проезде препятствия лучше всего:}

\begin{itemize}[nosep,leftmargin=2em]
\item [A.] Сохранять обычную посадку.
\item [B.] Приподняться над сиденьем.
\item [C.] Сжать бензобак коленями.
\item [D.] Максимально наклониться вперед.

\end{itemize}
\end{minipage}


\begin{minipage}{\linewidth}
\vspace{1em}\noindent\textbf{67. На двухполосной дороге с твердым покрытием песок и гравий чаще всего скапливаются:}

\begin{itemize}[nosep,leftmargin=2em]
\item [A.] В левой колее.
\item [B.] Ближе к центру дороги.
\item [C.] У краев дороги (у обочин).
\item [D.] По центру полосы.

\end{itemize}
\end{minipage}


\begin{minipage}{\linewidth}
\vspace{1em}\noindent\textbf{68. На мокрой дороге наиболее скользко обычно:}

\begin{itemize}[nosep,leftmargin=2em]
\item [A.] На мелком гравии.
\item [B.] В левой колее.
\item [C.] В правой колее.
\item [D.] В центральной части полосы.

\end{itemize}
\end{minipage}


\begin{minipage}{\linewidth}
\vspace{1em}\noindent\textbf{69. Чтобы безопасно проехать скользкий участок дороги, лучше всего:}

\begin{itemize}[nosep,leftmargin=2em]
\item [A.] Легко притормаживать задним тормозом во время движения по нему.
\item [B.] Снизить скорость до въезда на него.
\item [C.] Постепенно ускоряться.
\item [D.] Переключиться на пониженную передачу при въезде на него.

\end{itemize}
\end{minipage}


\begin{minipage}{\linewidth}
\vspace{1em}\noindent\textbf{70. Для остановки на мокром асфальте:}

\begin{itemize}[nosep,leftmargin=2em]
\item [A.] Использовать передний тормоз только если колесо начинает скользить.
\item [B.] Двигаться накатом до полной остановки.
\item [C.] Использовать только задний тормоз.
\item [D.] Использовать передний и задний тормоза одновременно.

\end{itemize}
\end{minipage}


\begin{minipage}{\linewidth}
\vspace{1em}\noindent\textbf{71. При движении по дорогам с твердым покрытием в местах скопления песка и гравия следует:}

\begin{itemize}[nosep,leftmargin=2em]
\item [A.] Переключиться на пониженную передачу и добавить газ.
\item [B.] Выжать сцепление.
\item [C.] Использовать только задний тормоз для замедления.
\item [D.] Избегать резких изменений скорости или направления движения.

\end{itemize}
\end{minipage}


\begin{minipage}{\linewidth}
\vspace{1em}\noindent\textbf{72. Если спущено заднее колесо:}

\begin{itemize}[nosep,leftmargin=2em]
\item [A.] Вы не сможете сохранить контроль.
\item [B.] Вы потеряете передачу мощности на заднее колесо.
\item [C.] Задняя часть мотоцикла будет вилять из стороны в сторону.
\item [D.] Рулевое управление станет «тяжелым».

\end{itemize}
\end{minipage}


\begin{minipage}{\linewidth}
\vspace{1em}\noindent\textbf{73. Если во время движения у вас пробило колесо, крепко держитесь за руль и:}

\begin{itemize}[nosep,leftmargin=2em]
\item [A.] Плавно задействуйте оба тормоза.
\item [B.] Используйте тормоз исправного колеса или избегайте торможения.
\item [C.] Продолжайте движение до ближайшей мастерской.
\item [D.] Используйте только тормоз спущенного колеса.

\end{itemize}
\end{minipage}


\begin{minipage}{\linewidth}
\vspace{1em}\noindent\textbf{74. Если во время движения спущено переднее колесо:}

\begin{itemize}[nosep,leftmargin=2em]
\item [A.] Используйте передний тормоз.
\item [B.] Вы не сможете управлять мотоциклом.
\item [C.] Используйте задний тормоз и съедьте к краю дороги.
\item [D.] Поверните вправо и примените оба тормоза.

\end{itemize}
\end{minipage}


\begin{minipage}{\linewidth}
\vspace{1em}\noindent\textbf{75. При спущенном переднем колесе:}

\begin{itemize}[nosep,leftmargin=2em]
\item [A.] Рулевое управление станет тяжелым.
\item [B.] Это никак не повлияет на управляемость.
\item [C.] Задняя часть мотоцикла будет вилять из стороны в сторону.
\item [D.] Вы потеряете эффективность тормозов.

\end{itemize}
\end{minipage}


\begin{minipage}{\linewidth}
\vspace{1em}\noindent\textbf{76. Если заклинило ручку газа:}

\begin{itemize}[nosep,leftmargin=2em]
\item [A.] Выжмите сцепление и используйте аварийный выключатель двигателя.
\item [B.] Переключитесь на нейтраль и примените оба тормоза.
\item [C.] Переключитесь на пониженную передачу и примените оба тормоза.
\item [D.] Используйте передний тормоз.

\end{itemize}
\end{minipage}


\begin{minipage}{\linewidth}
\vspace{1em}\noindent\textbf{77. На рисунке выше обычно лучше всего:}

\vspace{1ex}
\centering
\includegraphics[width=0.8\linewidth]{ q_77.jpeg }
\captionof*{figure}{ Иллюстрация полосы движения с тремя пронумерованными частями: 1 — левая, 2 — центральная, 3 — правая. }

\begin{itemize}[nosep,leftmargin=2em]
\item [A.] Увеличить скорость, чтобы быстро проехать.
\item [B.] Двигаться в позиции 3.
\item [C.] Двигаться в позиции 1.
\item [D.] Двигаться в позиции 2 или 3.

\end{itemize}
\end{minipage}


\begin{minipage}{\linewidth}
\vspace{1em}\noindent\textbf{78. На этом рисунке A, B и C едут в группе. Для обгона мотоциклист B должен:}

\vspace{1ex}
\centering
\includegraphics[width=0.8\linewidth]{ q_78.jpeg }
\captionof*{figure}{ Иллюстрация группы мотоциклов, обгоняющих автомобиль: A уже завершил обгон, B совершает обгон, C ожидает. }

\begin{itemize}[nosep,leftmargin=2em]
\item [A.] Завершить обгон перед мотоциклистом A.
\item [B.] Совершать обгон одновременно с C.
\item [C.] Вернуться в полосу и ехать рядом с A.
\item [D.] Вернуться в правую часть полосы позади A.

\end{itemize}
\end{minipage}


\begin{minipage}{\linewidth}
\vspace{1em}\noindent\textbf{79. Чтобы создать больше пространства в ситуации, изображенной выше, вам следует:}

\vspace{1ex}
\centering
\includegraphics[width=0.8\linewidth]{ q_79.jpeg }
\captionof*{figure}{ Мотоцикл едет по правой полосе шоссе, автомобиль перестраивается в правую полосу перед мотоциклом. }

\begin{itemize}[nosep,leftmargin=2em]
\item [A.] Ехать в правой части полосы.
\item [B.] Скорректировать скорость и оставаться по центру полосы.
\item [C.] Перестроиться в левую полосу.
\item [D.] Обогнать автомобиль слева от вас.

\end{itemize}
\end{minipage}


\begin{minipage}{\linewidth}
\vspace{1em}\noindent\textbf{80. Большинство аварий с участием мотоцикла и автомобиля:}

\begin{itemize}[nosep,leftmargin=2em]
\item [A.] Происходят из-за того, что водитель автомобиля выезжает на траекторию движения мотоцикла.
\item [B.] Приводят к незначительным травмам.
\item [C.] Происходят из-за того, что мотоциклист поворачивает налево перед транспортным средством.
\item [D.] Происходят из-за превышения скорости.

\end{itemize}
\end{minipage}


\begin{minipage}{\linewidth}
\vspace{1em}\noindent\textbf{81. Вы собираетесь перестроиться в левую полосу, как показано на рисунке выше. Чтобы заметить автомобили, проезжающие рядом с вами, лучше всего:}

\vspace{1ex}
\centering
\includegraphics[width=0.8\linewidth]{ q_81.jpeg }
\captionof*{figure}{ Мотоцикл по центру средней полосы, готовится к перестроению в левую полосу. }

\begin{itemize}[nosep,leftmargin=2em]
\item [A.] Замедлиться и посмотреть в левое зеркало.
\item [B.] Посмотреть в левое зеркало.
\item [C.] Ускориться и посмотреть в левое зеркало.
\item [D.] Повернуть голову влево.

\end{itemize}
\end{minipage}


\begin{minipage}{\linewidth}
\vspace{1em}\noindent\textbf{82. Чтобы подстроиться под дорожное движение и избежать мусора на повороте, показанном на рисунке, следует ехать:}

\vspace{1ex}
\centering
\includegraphics[width=0.8\linewidth]{ q_82.jpeg }
\captionof*{figure}{ Поворот дороги без разделительного барьера, где правая сторона имеет пронумерованные части: 1 — левая часть, 2 — центральная часть, 3 — правая часть. }

\begin{itemize}[nosep,leftmargin=2em]
\item [A.] В позиции 3.
\item [B.] В позиции 2.
\item [C.] Без наклона.
\item [D.] В позиции 1.

\end{itemize}
\end{minipage}


\begin{minipage}{\linewidth}
\vspace{1em}\noindent\textbf{83. Чтобы пересечь рельсы, идущие параллельно вашему курсу:}

\begin{itemize}[nosep,leftmargin=2em]
\item [A.] Постепенно переезжать их под острым углом.
\item [B.] Выехать на встречную полосу, чтобы пересечь их под углом 90 градусов.
\item [C.] Подождать, пока рельсы закончатся.
\item [D.] Отъехать в сторону, затем повернуть и пересечь их под углом не менее 45 градусов.

\end{itemize}
\end{minipage}


\begin{minipage}{\linewidth}
\vspace{1em}\noindent\textbf{84. На этом рисунке мотоциклист, находящийся в наиболее опасном положении, это:}

\vspace{1ex}
\centering
\includegraphics[width=0.8\linewidth]{ q_84.jpeg }
\captionof*{figure}{ Иллюстрация группы мотоциклистов, обгоняющих автомобиль: A уже завершил обгон, B совершает обгон, C ожидает. }

\begin{itemize}[nosep,leftmargin=2em]
\item [A.] Мотоциклист B.
\item [B.] Мотоциклист C.
\item [C.] Мотоциклист A.
\item [D.] Никто.

\end{itemize}
\end{minipage}


